\itemhead{2.3 例题 · 例题9（教材p.57）}{XJ-B1-C02-EX-P057-Y0534}
\label{XJ-B1-C02-EX-P057-Y0534}
{\footnotesize 湘教版教材p.57 / PDF第64页 / 2.3 例题 / 例题9}
\field{原题}{某化学试剂生产厂以 $x\,\mathrm{kg/h}$ 的速度运输生产某种产品（生产条件要求边生产边运输，且 $1\le x\le10$），每小时可获得利润 $100\left(5x+1-\dfrac3x\right)$ 元。
\subq{1}要使运输生产该产品 $2\,\mathrm h$ 获得的利润不低于 3000 元，求 $x$ 的取值范围。
\subq{2}要使运输生产 $900\,\mathrm{kg}$ 该产品获得的利润最大，该工厂应该选取何种运输生产速度？并求最大利润。}
\field{参考答案}{（1）$5x+1-3/x\ge15$ 得 $5x^2-14x-3\ge0$，结合 $1\le x\le10$ 得 $x\in[3,10]$。（2）总利润 $90000(5+1/x-3/x^2)=457500-270000(1/x-1/6)^2$，在 $x=6$ kg/h 时最大457500元。}
{\small \field{关联知识}{I03 基本不等式、等号与最值建模（教材帮p.52、53、54 / PDF61、62、63）；I04 一元二次不等式与参数分类（教材帮p.64、65、66 / PDF73、74、75）；F03 函数求值、解析式与函数建模（教材帮p.80、87、88、89、90、131 / PDF89、96、97、98、99、140）}
\field{方法与易错点}{列正数条件与固定量，构造基本不等式，解等号条件并代回实际限制。 易错：把一个下界误当最小值；多个等号条件不能同时实现；漏单位和整数限制。}
\field{答案依据}{例题原解位置：湘教版教材p.57及题库原图；独立编写可编辑参考解答。答案册此部分主要收练习、习题，不伪造例题答案册位置。}}
